\section{Discussion}
\label{sec:discussion}

\paragraph{RQ1: controlled correctness.}
The correctness and invariance gate passed without exceptions.  Cost equality
alone would allow subtle scheduling differences to remain hidden.  Identical
full-landmark expansion digests additionally show that eager, lazy, and staged
conditions performed the same successor-enumerating search under the frozen
tie breaker.  Their differences in time and counters can therefore be
localized to evaluation timing and OPEN promotions rather than to different
expanded state sets.

\paragraph{RQ2: staging saved real work.}
The prefix was not merely ceremonial.  Staging reduced median pivot and read
work by about 3\%, and much more on maze queries.  In operational terms,
\LazyFull{} evaluates all 32 terms whenever a cheap state becomes competitive;
\Staged{} spends four reads first and avoids the other 28 when that prefix
pushes the state below competitors.  The saving is exact and deterministic.
It is also heterogeneous and often zero, so the overall median understates
large savings in one regime and overstates the effect in another.

\paragraph{RQ3: saved work did not imply universal speed.}
Every intermediate evaluation requires another pop and requeue.  The staged
schedule averaged roughly 178 more requeues and 178 more pops per query than
ordinary lazy evaluation.  Its primary paired ratio of 1.0543 therefore points toward
a modest slowdown at the pooled median, while the 95\% interval includes
parity.  We interpret this as bounded negative evidence for universal use of
this fixed prefix, not as proof that staging can never help.  Indeed, the maze
maps show the expected positive regime.

The strong negative Spearman association between read saving and log time
ratio supports the proposed mechanism descriptively: queries that avoided a
larger fraction of reads tended to make staging relatively faster.  It cannot
establish that removing a particular read would cause the measured speedup.
Both variables co-vary with topology, query length, absolute number of
competitive nodes, cache locality, and Python priority-queue costs.

\paragraph{RQ4: topology and workload matter.}
Warehouses are the clearest failure regime.  The four-pivot prefix had zero
median saving on both maps, so it added a nearly pure heap cycle before the
same suffix work.  Mazes are the clearest success regime: the prefix filtered
enough candidates to save about 30\% of reads on both maps, outweighing
promotion overhead.  Random and room maps occupy intermediate or mixed
regimes.  This crossover is more informative than a single leaderboard rank:
it specifies when the simple schedule's mechanism is active.

Landmark construction changes the deployment decision at a different level.
If a map receives only a few queries, even a fast landmark search cannot repay
dozens of BFS preprocessing passes.  At high query volume the fixed cost is
small per query, and all landmark modes beat Manhattan under the study's
equal-map average.  The analysis deliberately charges the same build cost to
eager, lazy, and staged modes because they use the same table; staging is a
query-time schedule, not a preprocessing improvement.

\paragraph{Relation to prior work.}
The results instantiate the tradeoff identified by Lazy \Astar{}: heuristic
calls saved on surplus nodes versus extra OPEN manipulation
\citep{tolpin2013towards}.  They do not imply an optimal deployment policy and
do not challenge rational or learned methods, which may choose conditionally
whether to evaluate \citep{karpas2018rational,domshlak2012online}.  A natural
extension would use predecision features to choose between ordinary lazy and
staged schedules.  That is intentionally future work: fitting such a selector
after seeing these maps would change the question and require a new split.
